% Academic Research Paper Template
% For IEEE-style publications and academic research

\documentclass[conference]{IEEEtran}

% Packages
\usepackage{cite}
\usepackage{amsmath,amssymb,amsfonts}
\usepackage{algorithmic}
\usepackage{graphicx}
\usepackage{textcomp}
\usepackage{xcolor}
\usepackage{hyperref}

% Hyperref setup
\hypersetup{
    colorlinks=true,
    linkcolor=blue,
    filecolor=magenta,
    urlcolor=cyan,
    citecolor=blue,
}

% Document metadata
\title{Your Research Paper Title\\
{\large Subtitle if Needed}}

\author{
    \IEEEauthorblockN{Scott Weeden}
    \IEEEauthorblockA{
        \textit{Department of Computer Science} \\
        \textit{Texas Tech University}\\
        Lubbock, Texas, USA \\
        email@ttu.edu
    }
}

\begin{document}

\maketitle

\begin{abstract}
This is where you write your abstract. The abstract should be a concise summary of your research, typically 150-250 words. It should include: (1) the problem or motivation, (2) the approach or methodology, (3) key results, and (4) main conclusions. Keep it focused and clear.
\end{abstract}

\begin{IEEEkeywords}
machine learning, algorithm analysis, computational complexity, your keywords here
\end{IEEEkeywords}

\section{Introduction}
Introduce your research problem and motivation. Explain why this work is important and what gap in knowledge or practice it addresses.

\subsection{Background}
Provide necessary background information for understanding your work.

\subsection{Contributions}
Clearly state your main contributions:
\begin{itemize}
    \item First contribution
    \item Second contribution
    \item Third contribution
\end{itemize}

\section{Related Work}
Discuss previous research related to your work. Compare and contrast with your approach.

\section{Methodology}
Describe your approach, algorithms, or experimental setup.

\subsection{Problem Formulation}
Formally define the problem. Use mathematical notation where appropriate:

\begin{equation}
f(x) = \sum_{i=1}^{n} w_i x_i
\end{equation}

\subsection{Algorithm}
Present your algorithm or approach.

\section{Experimental Results}
Present and analyze your results.

\subsection{Experimental Setup}
Describe your experimental environment, datasets, and evaluation metrics.

\subsection{Results and Analysis}
Present results with figures and tables.

\begin{figure}[htbp]
\centerline{\includegraphics[width=0.45\textwidth]{figure1.png}}
\caption{Your figure caption here.}
\label{fig:example}
\end{figure}

\begin{table}[htbp]
\caption{Results Comparison}
\begin{center}
\begin{tabular}{|c|c|c|}
\hline
\textbf{Method} & \textbf{Accuracy} & \textbf{Time (s)} \\
\hline
Baseline & 85.2\% & 12.3 \\
\hline
Our Method & \textbf{92.7\%} & \textbf{8.1} \\
\hline
\end{tabular}
\label{tab:results}
\end{center}
\end{table}

\section{Discussion}
Discuss the implications of your results, limitations, and potential improvements.

\section{Conclusion}
Summarize your work and contributions. Mention future research directions.

\section*{Acknowledgments}
Acknowledge funding sources, collaborators, and others who contributed.

\bibliographystyle{IEEEtran}
\bibliography{references}

\end{document}
